\uuid{Iyo4}
\titre{Distance}
\auteur{bodin ; hanani ; mzari}
\organisation{exo7}
\licence{CC BY-NC-SA 4.0}
\datecreate{2020-01-20}
\niveau{L1}
\module{Géométrie}
\chapitre{Géométrie de l'espace}
\sousChapitre{Distance}
\difficulte{3}
\theme{Distance}
\exo7qcmid{lille-1.geometrie-espace.030}

\contenu{
\texte{
Pour ces questions, l'espace est muni d'un repère orthonormé direct  $(O, \vec{i}, \vec{j}, \vec{k})$.
}

\begin{question}
\qclassification{141.05}
\qtags{motcle=distance}
Soit  $D$ la droite : $\left\{\begin{array}{ccl}x&=&-1+3t\\y&=&1\\z&=&-1-t, \quad t\in \Rr  \end{array}\right.$ et $P$ un plan contenant $D$ à une distance de $1$ de l'origine. Une équation cartésienne de $P$ est :
\begin{answers}
      \bad{$y=1$}

      \good{$y=1$ ou $4x+3y+12z+13=0$}

      \bad {$y=1$ ou $x=1$}

      \bad {$x=1$ ou $y=1$ ou $z=1$}

\end{answers}
\begin{explanations}
Une représentation cartésienne de $D$ est : $\left\{\begin{array}{ccl}y&=&1\\x+3z+4&=&0  \end{array}\right.$. $P$ est différent du plan $x+3z+4=0$ et $P$ contient $D$, une équation cartésienne de $P$ est donc de la forme : $(y-1)+\alpha(x+3z+4)=0$. On calcule $\alpha$ de sorte que la distance de $P$ à l'origine soit égale à $1$.
\end{explanations}
\end{question}
}
